% Part: sets-functions-relations
% Chapter: relations
% Section: equivalence-relations

\documentclass[../../../include/open-logic-section]{subfiles}

\begin{document}

\olfileid{sfr}{rel}{eqv}

\olsection{तुल्यता संबंध}

एखाद्या संचावरील एकरूपता संबंध परावर्ती, सममित आणि संक्रमक असतो.
हे तिन्ही गुणधर्म असणारे $R$ संबंध वारंवार आढळतात.

\begin{defn}[तुल्यता संबंध]
परावर्ती, सममित आणि संक्रमक असणाऱ्या $R \subseteq A^2$ या संबंधाला
\emph{तुल्यता संबंध} म्हणतात. $A$ मधील !!^{element}s $x$ आणि $y$ यांना
\emph{$R$-तुल्य} म्हणतात, जर $Rxy$ असते.
\end{defn}

तुल्यता संबंधांमुळे \emph{तुल्यतावर्ग} ही संकल्पना मिळते. तुल्यता संबंध
प्रांताचे वेगवेगळ्या भागांत ``तुकडे'' करतो. प्रत्येक भागात सर्व वस्तू
एकमेकींशी संबंधित असतात; भिन्न भागांतील कोणत्याही वस्तू एकमेकींशी संबंधित
नसतात. काही वेळा या भागांबद्दल \emph{थेट} बोलणे उपयोगी ठरते. त्यासाठी
पुढील व्याख्या करू:

\begin{defn}\ollabel{def:equivalenceclass}
$R \subseteq A^2$ हा तुल्यता संबंध असू द्या. प्रत्येक $x \in A$ साठी,
$A$ मधील $x$ चा \emph{तुल्यतावर्ग} म्हणजे
$\equivrep{x}{R} = \Setabs{y \in A}{Rxy}$ हा संच. $R$ नुसार $A$ चा
\emph{भागाकारसंच} म्हणजे $\equivclass{A}{R} = \Setabs{\equivrep{x}{R}}{x \in A}$,
अर्थात या तुल्यतावर्गांचा संच.
\end{defn}

पुढील निष्कर्ष तुल्यतावर्गाच्या व्याख्येला आधार देतो: तुल्यतावर्ग हे खरोखर
$A$ चे विभाजन करणारे भाग असतात, हे तो सिद्ध करतो.

\begin{prop}
$R \subseteq A^2$ हा तुल्यता संबंध असेल, तर $Rxy$ तेव्हा आणि केवळ
तेव्हाच, जेव्हा $\equivrep{x}{R} = \equivrep{y}{R}$.
\end{prop}

\begin{proof}
डावीकडून उजवीकडे जाणाऱ्या दिशेसाठी, $Rxy$ असे समजा आणि
$z \in \equivrep{x}{R}$ असू द्या. व्याख्येनुसार $Rxz$. $R$ हा तुल्यता
संबंध असल्याने $Ryz$. (सविस्तर सांगायचे तर: $Rxy$ असल्याने आणि $R$
सममित असल्याने $Ryx$; तसेच $Rxz$ असल्याने आणि $R$ संक्रमक असल्याने
$Ryz$.) म्हणून $z \in \equivrep{y}{R}$. हे कोणत्याही अशा घटकासाठी
लागू असल्याने $\equivrep{x}{R} \subseteq \equivrep{y}{R}$. अगदी
त्याचप्रमाणे $\equivrep{y}{R} \subseteq \equivrep{x}{R}$. म्हणून
विस्तारात्मकतेनुसार $\equivrep{x}{R} = \equivrep{y}{R}$.

उजवीकडून डावीकडे जाणाऱ्या दिशेसाठी,
$\equivrep{x}{R} = \equivrep{y}{R}$ असे समजा. $R$ परावर्ती असल्याने
$Ryy$, म्हणून $y \in \equivrep{y}{R}$. मग
$\equivrep{x}{R} = \equivrep{y}{R}$ या गृहीतकामुळे
$y \in \equivrep{x}{R}$ सुद्धा. म्हणून $Rxy$.
\end{proof}

\begin{ex}
तुल्यता संबंधांचे एक चांगले उदाहरण शेषांक अंकगणितातून मिळते. कोणत्याही
$a$, $b$ आणि $n \in \PosInt$ साठी, $a$ ला $n$ ने भागल्यावर येणारी
बाकी आणि $b$ ला $n$ ने भागल्यावर येणारी बाकी सारखी असेल, तेव्हा आणि
केवळ तेव्हाच $a \equiv_n b$ असे म्हणू. (चिन्हांत मांडायचे तर:
$a \equiv_n b$ तेव्हा आणि केवळ तेव्हाच, जेव्हा एखाद्या $k \in \Int$
साठी $a - b = kn$.) मग कोणत्याही $n$ साठी $\equiv_n$ हा तुल्यता
संबंध असतो. $\equiv_n$ मुळे नेमके $n$ भिन्न तुल्यतावर्ग मिळतात; म्हणजे
$\equivclass{\Nat}{\equiv_n}$ मध्ये $n$ !!{element}s असतात. ते असे:
$n$ ने निःशेष भाग जाणाऱ्या संख्यांचा संच, म्हणजे $\equivrep{0}{\equiv_n}$;
$n$ ने भागल्यावर बाकी $1$ येणाऱ्या संख्यांचा संच, म्हणजे
$\equivrep{1}{\equiv_n}$; \ldots; आणि $n$ ने भागल्यावर बाकी $n-1$
येणाऱ्या संख्यांचा संच, म्हणजे $\equivrep{n-1}{\equiv_n}$.
\end{ex}

\begin{prob}
कोणत्याही $n \in \PosInt$ साठी $\equiv_n$ हा तुल्यता संबंध आहे आणि
$\equivclass{\Nat}{\equiv_n}$ मध्ये नेमके $n$ सदस्य आहेत, हे दाखवा.
\end{prob}

\end{document}
